\documentclass[11pt]{report}
\usepackage[utf8]{inputenc}

\usepackage{graphicx}
\usepackage{geometry}
\usepackage{enumitem}
\usepackage{hyperref}
\usepackage{colortbl}
\usepackage{array}
\usepackage{float}
\usepackage{tabularx}
\usepackage{longtable}

\hypersetup{
    colorlinks=true,
    linkcolor=black,
    filecolor=magenta,      
    urlcolor=cyan,
    pdftitle={Requirement Analysis and Specification Document},}

\graphicspath{ {Images/} }
\geometry{a4paper, margin=100px}

\title{CodeKataBattle - Requirement Analysis and Specification Document}
% https://github.com/FilippoRiva/RivaSulfaro
\author{Filippo Riva, Antonio Sulfaro}
\date{\today}

\begin{document}

\begin{titlepage}
    \begin{figure}[t]
        \centering
        %\includegraphics[width=250px]{/PolimiLogo}
    \end{figure}
    \begin{centering}
       \maketitle
    \end{centering}
\end{titlepage}

\tableofcontents

%  1  --------------------------------------------------------------------------------------------%

\chapter{Introduction}
    Challenges may be a good way to give students the opportunity and the willpower to learn,
    understand and study some topics, to compare themselves with others and to enjoy more
    what they study.

    In IT this statement is true more then ever since programming is a complex challenge,
    especially in working environments.
    CKB makes the challenge real, giving to educators the possibility to teach 
    and test their students through real projects, while giving to students the possibility to have fun 
    and improve their knowledge of the programming languages as if it was a game.

    The entire platform can be used by students to improve their programming skills, by schools 
    to make internal tournaments, by companies to test their employees and by many other stakeholders. 
    \\


    \section{Purpose}
    The CodeKataBattle project, as it's name might suggest, is taking inspiration from the concept 
    of a kata.
    A kata is an exercise in karate where you repeat a form many, many times, making
    little improvements each time you repeat the movement.
    The main idea is to apply this concept to competitive programming, giving students many
    problems to solve and making them learn by trial and error.

    Let's now introduce the main goals of the CKB project.
    \subsubsection{G1 - Allow educators to create Kata Tournaments}
    \subsubsection{G2 - Allow educators to create Kata Battles for a specific Kata Tournament} 
    \subsubsection{G3 - Allow educators to evaluate students automatically via the system and optionally via a manual scoring}
    \subsubsection{G4 - Allow students to subscribe to Kata Tournaments available on the CKB platform}
    \subsubsection{G5 - Allow students to join Kata Battles of a certain Kata Tournament alone or as a team}
    \subsubsection{G6 - Allow students to compete in specific Kata Tournament}

    \section{Scope}
        \subsection{Phenomena}
            \subsubsection{World Phenomena}
                \begin{itemize}
                    \item [$W01$] Educators create/prepare new katas.
                    \item [$W02$] Students programs the solutions to katas.
                    \item [$W03$] Students fork CKB's official GitHub repositories.
                    \item [$W04$] Students upload their solutions to forks of CKB's official GitHub repositories.
                \end{itemize}
            \subsubsection{World Controlled Shared Phenomena}
                \begin{itemize}
                    \item [$S01$] Students sign up to the system.
                    \item [$S02$] Educators sign up to the system.
                    \item [$S03$] Educators create Kata Tournaments.
                    \item [$S04$] Educators create Kata Battles.
                    \item [$S05$] Educators evaluate students' solutions.
                    \item [$S06$] Educators who created a Kata Tournament grant to other educators the permission to create battles within that tournament.
                    \item [$S07$] Students subscribe to a Kata Tournament.
                    \item [$S08$] Students join a Kata Battle.
                    \item [$S09$] Students invite other students to participate to Kata Battles as a group.
                \end{itemize}
            \subsubsection{Machine Controlled Shared Phenomena}
                \begin{itemize}
                    \item [$M01$] The system notifies users about the beginning of a new tournament
                    \item [$M02$] The system notifies users about the start of a kata battle. 
                    \item [$M03$] The system notifies users about the final ranks of a kata battle.
                    \item [$M04$] The system notifies users about their subscription to a kata battle.
                    \item [$M05$] The system pulls new code from GitHub forks and updates the scores accordingly
                \end{itemize}
    \section{Definitions, Acronyms, Abbreviations}
    \subsection{Definitions}
    \emph{Team} $\rightarrow$ A group of one or more people taking part to a kata battle
    \\ \\
    \emph{GitHub} $\rightarrow$ 
    \\ \\
    \emph{GitHub repositories} $\rightarrow$ 
    \\ \\
    \emph{Consolidation Stage} $\rightarrow$ 
    \\ \\
    \emph{Tournament Rank} $\rightarrow$ 
    \\ \\
    \emph{Battle Rank} $\rightarrow$ 
    \\ \\
    \emph{Score} $\rightarrow$ 
    \\ \\
    \emph{Organizer} $\rightarrow$ 
    \\ \\
    \emph{} $\rightarrow$ 

    \subsection{Acronyms}
    \emph{CKB}  $\rightarrow$ the CodeKataBattle platform
    \\ \\
    \emph{KB}   $\rightarrow$ a Kata Battle.
    \\ \\
    \emph{KT}   $\rightarrow$ a Kata Tournament.
    \\ \\
    \emph{WPx}  $\rightarrow$ x-th World Phenomena.
    \\ \\
    \emph{SPx}  $\rightarrow$ x-th Shared Phenomena.
    \\ \\
    \emph{Rx}   $\rightarrow$ x-th Requirement.
    \\ \\
    \emph{Dx}   $\rightarrow$ x-th Domain Assumption.
    \\ \\
    \emph{Sx}   $\rightarrow$ x-th Scenario.
    \\ \\
    \emph{repo} $\rightarrow$ GitHub repository.

    \section{Revision History}
    \begin{itemize}
        \item November 29 , 2023: version 1.0 (document structure)
    \end{itemize}

    \section{Reference Documents}
    \emph{Reference book} $\rightarrow$ Hans Van Vilet - Software Engineering: Principles and Practice
    \\ \\
    \emph{Alloy official documentation} $\rightarrow$ https://alloytools.org/documentation.html
    \\ \\
    \section{Document Structure}
    \begin{itemize}
        \item \textbf{Section 1: Introduction}\\
        \item \textbf{Section 2: Overall Description}\\
        \item \textbf{Section 3: Specific Requirements}\\
        \item \textbf{Section 4: Formal Analysis through Alloy} \\
        \item \textbf{Section 5: Effort spent} \\
    \end{itemize}

%  2  --------------------------------------------------------------------------------------------%
\chapter{Overall Description}
\section{Product Perspective}

\subsection{Class diagram}

\subsection{Scenarios}
\subsubsection{S1 - Educator signs up}
\subsubsection{S2 - Educator creates a KT}
\subsubsection{S3 - Educator creates a KB}
\subsubsection{S4 - Educator manually evaluates the students' solutions to a KB}
\subsubsection{S5 - Student signs up}
\subsubsection{S6 - Student subscribes to a KT}
\subsubsection{S7 - Student joins a KB}
\subsubsection{S8 - A KB starts}
\subsubsection{S9 - Student pushes a new solution to his GitHub repo and the system tests it}
\subsubsection{S10 - A KB ends}

\section{Product Functions}
\begin{itemize}
    \item \textbf{Users registration and login}
    
    These functions will be available to both students and educators.

    \item \textbf{Create KTs and KBs}


    \item \textbf{Subscribe to a KT}
    \item \textbf{Join and compete ina KB alone or as a team}
    \item \textbf{Automatic Evaluation and GitHub integration}

    via submission API, GitHub Actions, quality level aspects\ldots description of what is evaluated

    \item \textbf{Notifications}
    \item \textbf{Public Rankings}
\end{itemize}

\section{User Characteristics}
\subsection{Educator}
The educator is an individual from a university, school, company
or organization that subscribes to the CKB platform in order to organize
competitions between students. This type of user can create KT about any
kind of topic they want and can create KB related to the KT topics.
The educators want to constantly see the tournament rankings and the current
scores assigned to each team participating to a KB of that tournament.
During the KT the educator might evaluate students’ solutions to the KB
problems if they think a manual evaluation is needed.
Educators are considered organizer when they create a tournament or are
granted the permission to create KB by a tournament creator (close
Tournament)
\subsection{Student}
The student is any kind of individual who wants to participate to a KT\@.


\section{Domain Assumptions}
    \begin{enumerate}[label=(D\arabic*), font=\itshape]
        \item Educators and students have an internet connection
        \item Students know how to use GitHub and fork other repositories
        \item Students know how to program GitHub Actions to notify the CKB platform
        \item Educators create KT and KB properly.
    \end{enumerate}


%  3  --------------------------------------------------------------------------------------------%
\chapter{Specific Requirements}

\section{External Interface Requirements}

\subsection{User Interfaces}
The website interface, clean easy to use

\subsection{Hardware Interface}
A computer with an installed browser since all functionalities are being offered through a website

\subsection{Software Interfaces}
The submission API to be used from GitHub Actions to automate the pull mechanism.

\subsection{Communication Interfaces}
The submission API to be used from GitHub Actions to automate the pull mechanism.


\section{Functional Requirements}
Requirements : 
\begin{table}[]
\begin{tabular}{|l|p{130mm}|}
\hline
R1 & The system must allow educators and students to sign up. \\ \hline
R2 & The system must allow educators and students to identify themselves via assign in procedure. \\ \hline
R3 & The system must allow educators to create a new KT. \\ \hline
R4 & The system must allow educators who are creating a KT to choose a KT name and description. \\ \hline
R5 & The system must allow educators who are creating a KT to insert a KT subscription deadline. \\ \hline
R6 & The system must allow educators who created a KT to grant the permission to other educators to become organizers of that KT. \\ \hline
R7 & The system must notify all students about the creation of a new KT. \\ \hline
R8 & The system must allow students to subscribe to a KT before its subscription deadline. \\ \hline
R9 & The system must allow students and educators to read the descriptions of every KT. \\ \hline
R10 & The system must allow students and educators to read the description of all the KB of every KT. \\ \hline
R11 & The system must allow students, before the start of a KB, to invite other students to join their team if the team size doesn't exceed the maximum one defined by the organizer. \\ \hline
R12 & The system, when the registration deadline of a KB expires, must create a new GitHub repository containing description, test cases and build automations scripts of the corresponding KB. \\ \hline
R13 & The system must notify all the students about the start of a KB that they have previously joined. \\ \hline
R14 & The system must provide a submission API to be notified about a change of the source code of a specific team. \\ \hline
R15 & The system must obtain the new source code of a certain team every time it is notified via its submission API. \\ \hline
R16 & The system must automatically run static analysis tools over a team’s code each time it receives a new update and update the corresponding team’s score. \\ \hline
R17 & The system must automatically run tests over a team’s code each time it receives a new update and update the corresponding team’s score. \\ \hline
R18 & The system must update the KB rankings each time a team’s score change. \\ \hline
R19 & The system must allow students and educators to see the tournament rankings of every KT. \\ \hline
R20 & The system must allow organizers of a specific KT to create KBs. \\ \hline
R21 & The system must allow organizers who are creating a KB to specify its description. \\ \hline
R22 & The system must allow organizers who are creating a KB to specify its testcases. \\ \hline
R23 & The system must allow organizers who are creating a KB to specify its automation scripts. \\ \hline
R24 & The system must allow organizers who are creating a KB to set the minimum and maximum number of students per team for that KB. \\ \hline
% split table ---- 
\end{tabular}
\end{table}

\begin{table}[]
\begin{tabular}{|l|p{130mm}|}
\hline
% split table ---- 
R25 & The system must allow organizers who are creating a KB to set its registration deadline. \\ \hline
R26 & The system must allow organizers who are creating a KB to set its submission deadline. \\ \hline
R27 & The system must allow organizers who are creating a KB to set if a manual evaluation of the KB results is needed. \\ \hline
R28 & The system must allow organizers who are creating a KB to set which quality level aspects are to be automatically evaluated by the system. \\ \hline
R29 & The system must notify all students subscribed to a specific KT about the creation of a new KB in that tournament. \\ \hline
R30 & The system must allow educators and students to sign up \\ \hline
\end{tabular}
\end{table}

\subsection{Use case diagrams} 

\subsection{Sequence diagrams} 

\subsection{Activity Diagrams}

\section{Performance requirements}

\section{Design constraints}

\subsection{Standard compliance}

\subsection{Hardware Constraints}

\section{Software system attributes}

\subsection{Reliability}

\subsection{Availability}

\subsection{Security}

\subsection{Maintainability}

\subsection{Portability}


%  4  --------------------------------------------------------------------------------------------%
\chapter{Formal analysis using Alloy}

\section{Alloy Code}

\section{Alloy Worlds}
\subsection{First World}
\section{..}

%  5  --------------------------------------------------------------------------------------------%
\chapter{Effort spent}
\begin{table}[H]
\centering
\begin{tabular}{|l|l|}
\hline
\textbf{Student}       & \textbf{Effort} \\ \hline
Filippo Riva           & 8h             \\ \hline
Antonio Sulfaro        & 8h             \\ \hline
\end{tabular}
\end{table}

\end{document}